\section{API 仕様}

本章では \texttt{src/core/} が提供する公開 API と主要な型を示す。
UI 層はこれらの関数・型にのみ依存し、DOM 操作は core 内の
\texttt{rasterize.ts} に閉じ込められている。

\subsection{型定義（types.ts）}

\begin{apibox}{AppState}
\texttt{AppState} は UI 全体の単一情報源であり、永続化・共有の単位でもある。

\begin{lstlisting}
type AppState = {
  text: string;
  fontId: string;
  width: number;        // グリフ幅[px]
  height: number;       // グリフ高さ[px]
  size: number;         // 描画フォントサイズ[px]
  threshold: number;    // 2値化閾値 0-255
  offsetX: number;
  offsetY: number;
  bold: boolean;
  format: FormatOptions;
  overrides: Record<number, Matrix>; // charIndex→手動編集結果
};
\end{lstlisting}
\end{apibox}

\begin{apibox}{Glyph / Matrix}
\begin{lstlisting}
type Bit = 0 | 1;
type Matrix = Bit[][]; // matrix[row][col]
type Glyph = {
  char: string;
  width: number;
  height: number;
  matrix: Matrix;
};
\end{lstlisting}
\end{apibox}

\begin{apibox}{FormatOptions}
\begin{lstlisting}
type FormatOptions = {
  language: OutputLanguage;  // c|cpp|arduino|python|...
  structure: Structure;      // matrix2d|matrix3d|flat|bitpack-row|bitpack-col
  dataType: string;          // "uint8_t" 等
  bitOrder: BitOrder;        // "msb" | "lsb"
  radix: Radix;              // "bin" | "hex" | "dec"
  invert: boolean;
  variableName: string;
  itemsPerLine: number;
  includeCharComment: boolean;
  includeAsciiArt: boolean;
  padByte: boolean;
  indent: string;
  prefix: string;
  suffix: string;
};
\end{lstlisting}
\end{apibox}

\subsection{ラスタライズ（rasterize.ts）}

\begin{funcblock}{rasterizeChar(char: string, opts: RasterizeOptions): Glyph}
指定文字を Canvas に描画し、各ピクセルのアルファ値を
\texttt{opts.threshold} で 2 値化して \texttt{Glyph} を返す。

\apisection{引数}
\begin{itemize}[leftmargin=*,itemsep=1pt]
  \item \texttt{char}: 対象文字（1文字）。
  \item \texttt{opts.fontFamily}: CSS の \texttt{font-family}。
  \item \texttt{opts.width}, \texttt{opts.height}: 出力行列のサイズ。
  \item \texttt{opts.size}: 描画フォントの \texttt{px} サイズ。
  \item \texttt{opts.threshold}: $[0,255]$ のアルファ閾値。
  \item \texttt{opts.offsetX}, \texttt{opts.offsetY}: 描画オフセット。
  \item \texttt{opts.bold}: \texttt{true} で \texttt{bold} 書体。
\end{itemize}

\apisection{戻り値}
\texttt{Glyph}。\texttt{matrix[row][col]} は 0/1 の \texttt{Bit}。
\end{funcblock}

\subsection{エンコード（encode.ts）}

\begin{funcblock}{packRow(matrix: Matrix, bitOrder: BitOrder, padByte: boolean): number[]}
\texttt{matrix} を行方向に走査して 1bit/pixel にパックする。
\texttt{padByte=true} の場合、各行を 8bit 境界にゼロ埋めする。
\end{funcblock}

\begin{funcblock}{packColumn(matrix: Matrix, bitOrder: BitOrder, padByte: boolean): number[]}
列方向パック版。OLED/LCD の「縦8ドット=1バイト」フォーマットに相当。
\end{funcblock}

\begin{funcblock}{formatNumber(n: number, radix: Radix, bitWidth: number): string}
\texttt{n} を指定基数で文字列化する。\texttt{hex} は \texttt{0x} 接頭辞、
\texttt{bin} は \texttt{0b} 接頭辞、\texttt{dec} は素の 10 進。
\texttt{bitWidth} に応じて 0 埋めする。
\end{funcblock}

\subsection{フォーマット（format.ts）}

\begin{funcblock}{formatOutput(glyphs: Glyph[], opts: FormatOptions): string}
グリフ列とフォーマット設定から、ユーザーに表示する最終的な
出力文字列を生成する。言語別に以下のヘルパへディスパッチする。

\apisection{ディスパッチ先}
\begin{itemize}[leftmargin=*,itemsep=1pt]
  \item \texttt{formatC} (C / C++ / Arduino)
  \item \texttt{formatPython}
  \item \texttt{formatJS} (JavaScript / TypeScript)
  \item \texttt{formatRust}
  \item \texttt{formatGo}
  \item \texttt{formatJSON}
  \item \texttt{formatMarkdown}
  \item \texttt{formatSymbols} (記号アート、\texttt{■/□})
  \item \texttt{formatPlain} (生の 0/1 テキスト)
\end{itemize}
\end{funcblock}

\subsection{フォント管理（fonts.ts）}

\begin{funcblock}{loadFont(font: FontDef): Promise<void>}
\texttt{FontDef.url} が指す TTF/OTF を \texttt{FontFace} で読み込み、
\texttt{document.fonts} に登録する。同梱フォント用。
\end{funcblock}

\begin{funcblock}{registerCustomFont(file: File): Promise<FontDef>}
ユーザーがアップロードした \texttt{File} を \texttt{FontFace} 経由で
登録し、ID を付与した \texttt{FontDef} を返す。
以降の \texttt{fontId} 指定で利用可能になる。
\end{funcblock}

\subsection{共有リンク（share.ts）}

\begin{funcblock}{encodeStateToURL(state: AppState): string}
\texttt{overrides} を除く \texttt{AppState} を JSON 化し、base64url で
圧縮した文字列を \texttt{?s=} クエリに載せた URL を返す。
\end{funcblock}

\begin{funcblock}{decodeStateFromURL(): Partial<AppState> | null}
現在の URL から \texttt{?s=} を読み出し、復元した \texttt{AppState} を返す。
不正な値や欠損フィールドは無視し、欠けた箇所は既定値で補完される。
\end{funcblock}

\begin{notebox}[overrides を共有しない理由]
ピクセル編集結果（\texttt{overrides}）はサイズが大きく URL に載せづらいため、
共有リンクには含めない。共有先で再生成するか、
JSON エクスポート機能（\texttt{OutputPanel} のダウンロード）を用いる。
\end{notebox}

\subsection{React Hooks}

\begin{funcblock}{useAppState(initial: AppState): [AppState, Dispatch<Action>]}
\texttt{useReducer} を用いた状態管理。サポートするアクション:
\texttt{setText}, \texttt{setFontId}, \texttt{setSize},
\texttt{setWidth}, \texttt{setHeight}, \texttt{setFontSize},
\texttt{setThreshold}, \texttt{setOffsetX}, \texttt{setOffsetY},
\texttt{setBold}, \texttt{setFormat}, \texttt{setOverride},
\texttt{clearOverride}, \texttt{clearAllOverrides}, \texttt{replaceState}。
\end{funcblock}

\begin{funcblock}{useGlyphs(state: AppState): Glyph[]}
\texttt{text}, \texttt{fontId}, 幾何パラメータ、\texttt{overrides}
の変化にのみ反応して再計算するメモ化ラッパ。
\end{funcblock}

\begin{funcblock}{usePersist(state: AppState)}
\texttt{AppState} を \texttt{localStorage["font-to-bin.state.v1"]} に
自動保存する副作用フック。\texttt{overrides} は除外する。
\end{funcblock}

\newpage
